%!TEX TS-program = pdflatex

\documentclass[aspectratio=169,10pt]{beamer}

% -----------------------------------------------------------------------------
% Быстрая сборка: pdfLaTeX, без TikZ, PGFPlots и fontspec.
% В Overleaf выберите Compiler = pdfLaTeX.
% -----------------------------------------------------------------------------
\usepackage[T2A]{fontenc}
\usepackage[utf8]{inputenc}
\usepackage[english,russian]{babel}
\usepackage{cmap}
\usepackage{amsmath,amssymb,mathtools}
\usepackage{booktabs}
\usepackage{array}
\usepackage{etoolbox}
\usepackage{PTSans}
\usepackage{graphicx}
\graphicspath{{figures/}}

\renewcommand{\familydefault}{\sfdefault}
\renewcommand{\rmdefault}{\sfdefault}

\newbool{russian}
\booltrue{russian}
\usepackage{HSE-theme/beamerthemeHSE}

\hypersetup{bookmarks=false}

\setbeamertemplate{blocks}[rounded][shadow=false]
\setbeamerfont{frametitle}{size=\Large}
\setbeamerfont{framesubtitle}{size=\small}
\setbeamerfont{block title}{size=\normalsize}
\setbeamerfont{block body}{size=\small}

\newcommand{\R}{\mathbb{R}}
\newcommand{\norm}[1]{\left\lVert #1\right\rVert}
\newcommand{\ip}[2]{\left\langle #1,#2\right\rangle}
\newcommand{\grad}{\nabla}
\newcommand{\Hess}{\nabla^2}
\newcommand{\argmin}{\operatorname*{arg\,min}}

\newcommand{\coursebox}[2][2.65cm]{%
  \fbox{\parbox[c][0.92cm][c]{#1}{\centering #2}}%
}

\title[Методы оптимизации в ML]{Методы оптимизации\\в машинном обучении}
\subtitle{Лекция 2. Геометрия оптимизации: производная, градиент и гессиан}
\author[Н. Лукьяненко]{Никита Лукьяненко}
\institute{Совместная программа СМОЛ ГУ и ФКН НИУ ВШЭ}
\date{2026/27 учебный год}

\begin{document}

% 1 ---------------------------------------------------------------------------
\frame[plain]{\titlepage}

% 2 ---------------------------------------------------------------------------
\begin{frame}{От первой лекции --- к следующему вопросу}
\vspace{0.25em}

\begin{columns}[T,onlytextwidth]
\column{0.48\textwidth}
На первой лекции мы научились записывать обучение модели как задачу
\[
  x^\star\in\argmin_x f(x).
\]

Теперь представим, что мы находимся в некоторой текущей точке $x$.

\begin{alertblock}{Главный вопрос сегодня}
Как понять, \textbf{куда двигаться из $x$}, чтобы значение $f$ стало меньше?
\end{alertblock}

\column{0.48\textwidth}
\centering
\includegraphics[width=0.98\linewidth]{optimization_step_direction.pdf}
\end{columns}
\end{frame}

% 3 ---------------------------------------------------------------------------
\begin{frame}{Любой итерационный метод: одна простая схема}
\vspace{0.55em}

Во многих алгоритмах новая точка получается из старой по правилу
\[
  \boxed{x_{k+1}=x_k+h_k.}
\]

\begin{columns}[T,onlytextwidth]
\column{0.47\textwidth}
\begin{block}{Направление}
Куда должен смотреть вектор $h_k$?
\end{block}

\column{0.47\textwidth}
\begin{block}{Длина шага}
Насколько большим должен быть $\norm{h_k}$?
\end{block}
\end{columns}

\vspace{0.45em}

Сегодня почти полностью занимаемся \textbf{направлением}. Выбор длины шага будет отдельной темой курса.

\end{frame}

% 4 ---------------------------------------------------------------------------
\begin{frame}{Основная идея лекции}
\vspace{0.55em}

Сложную функцию трудно анализировать целиком. Поэтому около текущей точки $x$ мы заменяем её простой моделью.

\vspace{0.5em}
\centering
\coursebox[4.5cm]{\textbf{Линейная модель}\\первый порядок}
\qquad
$\Longrightarrow$
\qquad
\coursebox[4.0cm]{$\grad f(x)$\\градиент}

\vspace{0.65em}

\coursebox[4.5cm]{\textbf{Квадратичная модель}\\второй порядок}
\qquad
$\Longrightarrow$
\qquad
\coursebox[4.0cm]{$\Hess f(x)$\\гессиан}

\vspace{0.75em}

\begin{alertblock}{Одна история на всю лекцию}
\centering
Как поведение $f$ в маленькой окрестности точки $x$ подсказывает нам направление движения?
\end{alertblock}

\end{frame}

% 5 ---------------------------------------------------------------------------
\begin{frame}{Изменение параметров --- это вектор}
\small
\vspace{0.10em}

Пусть $x\in\R^d$ --- текущие параметры, а $h\in\R^d$ --- их изменение. Тогда
\[
  \boxed{x_{\mathrm{new}}=x+h.}
\]

\begin{columns}[T,onlytextwidth]
\column{0.46\textwidth}
Например,
\[
  x=\begin{pmatrix}2\\1\end{pmatrix},
  \qquad
  h=\begin{pmatrix}-1\\2\end{pmatrix}
\]
и поэтому
\[
  x+h=\begin{pmatrix}1\\3\end{pmatrix}.
\]

\begin{block}{Что кодирует $h$?}
Его направление отвечает на вопрос «куда двигаться?», а длина $\norm{h}$ --- «насколько далеко?».
\end{block}

\column{0.50\textwidth}
\centering
\includegraphics[width=0.76\linewidth]{vector_update.pdf}
\end{columns}
\end{frame}

% 6 ---------------------------------------------------------------------------
\begin{frame}{Длина изменения}
\vspace{0.55em}

Для евклидовой нормы
\[
  \norm{h}_2=\sqrt{h_1^2+\cdots+h_d^2}.
\]

\begin{block}{Геометрический смысл}
$\norm{h}_2$ --- обычная длина вектора: насколько далеко мы ушли в пространстве параметров.
\end{block}

Например, для
\[
 h=\begin{pmatrix}-1\\2\end{pmatrix}
\]
имеем
\[
 \norm{h}_2=\sqrt{5}.
\]

\vspace{0.2em}
\small
Дальше индекс $2$ будем часто опускать и писать просто $\norm{h}$.

\end{frame}

% 7 ---------------------------------------------------------------------------
\begin{frame}{Как сравнивать два направления?}
\small
\vspace{0.05em}

\begin{columns}[T,onlytextwidth]
\column{0.42\textwidth}
В $\R^d$ привычное евклидово скалярное произведение имеет вид
\[
  \boxed{x^\top v=\sum_{i=1}^d x_i v_i.}
\]

Его знак уже несёт геометрическую информацию:
\[
\begin{aligned}
  x^\top v_+&>0 && \text{--- острый угол},\\
  x^\top v_0&=0 && \text{--- прямой угол},\\
  x^\top v_-&<0 && \text{--- тупой угол}.
\end{aligned}
\]

\begin{block}{Геометрическая идея}
Скалярное произведение позволяет числом измерить, насколько два направления согласованы друг с другом.
\end{block}

\column{0.54\textwidth}
\centering
\includegraphics[width=0.98\linewidth]{dot_product_3d_signs.pdf}
\end{columns}

\end{frame}

% 8 ---------------------------------------------------------------------------
\begin{frame}{Самопроверка: что такое скалярное произведение?}
\vspace{0.75em}

\begin{alertblock}{Вопрос}
Пусть $V$ --- вещественное векторное пространство. Каким свойствам должно удовлетворять отображение
\[
  \langle\cdot,\cdot\rangle:V\times V\to\R,
\]
чтобы называться \textbf{скалярным произведением}?
\end{alertblock}

\vspace{0.65em}

\centering
\Large
Попробуйте вспомнить свойства до следующего слайда.

\vspace{0.55em}
\normalsize
Подсказка: нужны свойства, связывающие скалярное произведение с линейной структурой пространства и с понятием длины.

\end{frame}

% 9 ---------------------------------------------------------------------------
\begin{frame}{Самопроверка: определение скалярного произведения}
\small

\begin{block}{Вопрос}
Каким свойствам должно удовлетворять $\langle\cdot,\cdot\rangle:V\times V\to\R$?
\end{block}

\begin{block}{Определение}
Отображение $\langle\cdot,\cdot\rangle$ называется \textbf{скалярным произведением}, если для любых $u,v,w\in V$ и $\alpha,\beta\in\R$ выполнены:
\[
\begin{aligned}
\text{билинейность:}\quad
&\langle \alpha u+\beta w,v\rangle
 =\alpha\langle u,v\rangle+\beta\langle w,v\rangle,\\[-0.05em]
&\langle u,\alpha v+\beta w\rangle
 =\alpha\langle u,v\rangle+\beta\langle u,w\rangle;\\[0.25em]
\text{симметрия:}\quad
&\langle u,v\rangle=\langle v,u\rangle;\\[0.25em]
\text{положительная определённость:}\quad
&\langle u,u\rangle\ge0,
\qquad
\langle u,u\rangle=0\Longleftrightarrow u=0.
\end{aligned}
\]
\end{block}

Для стандартного скалярного произведения в $\R^d$:
\[
  \langle u,v\rangle=u^\top v=\sum_{i=1}^d u_i v_i.
\]

\end{frame}

% 8 ---------------------------------------------------------------------------
\begin{frame}{Скалярное произведение и угол}
\vspace{0.30em}

Для ненулевых векторов
\[
  \boxed{u^\top v=\norm{u}\,\norm{v}\cos\theta,}
\]
где $\theta\in[0,\pi]$ --- меньший угол между направлениями $u$ и $v$.

\begin{columns}[T,onlytextwidth]
\column{0.45\textwidth}
\vspace{0.30em}
\begin{itemize}
  \item $\theta<90^\circ$: $u^\top v>0$;
  \item $\theta=90^\circ$: $u^\top v=0$;
  \item $\theta>90^\circ$: $u^\top v<0$.
\end{itemize}

\small
Именно зависимость от угла позже позволит понять, почему движение против градиента уменьшает функцию.

\column{0.51\textwidth}
\centering
\includegraphics[width=0.95\linewidth]{angle_dot_product.pdf}
\end{columns}

\end{frame}


% 9 ---------------------------------------------------------------------------
\begin{frame}{Матрица как преобразование пространства}
\vspace{0.25em}

Пусть
\[
A=\begin{pmatrix}2&0\\0&1/2\end{pmatrix}.
\]
Тогда для любого $x=(x_1,x_2)^\top$
\[
  Ax=
  \begin{pmatrix}
    2x_1\\[0.15em]
    x_2/2
  \end{pmatrix}.
\]

\begin{columns}[T,onlytextwidth]
\column{0.45\textwidth}
\begin{block}{Что происходит геометрически?}
Первая координата каждой точки удваивается, а вторая уменьшается вдвое. Поэтому единичная окружность превращается в эллипс.
\end{block}

\small
Матрица действует не только на один нарисованный вектор: она одновременно преобразует \textbf{всё пространство}.

\column{0.51\textwidth}
\centering
\includegraphics[width=0.88\linewidth]{matrix_transform.pdf}
\end{columns}

\end{frame}


% 10 --------------------------------------------------------------------------
\begin{frame}{Квадратичная форма: матрица превращается в число}
\vspace{0.5em}

Если $x\in\R^d$ и $A\in\R^{d\times d}$, то
\[
  x^\top A x\in\R.
\]

\begin{block}{Следите за размерностями}
\[
  x\in\R^d,
  \qquad
  Ax\in\R^d,
  \qquad
  x^\top(Ax)\in\R.
\]
\end{block}

Например,
\[
 A=\begin{pmatrix}2&0\\0&5\end{pmatrix}
 \quad\Longrightarrow\quad
 x^\top A x=2x_1^2+5x_2^2.
\]

\small
Позже именно такие выражения будут описывать кривизну функции.

\end{frame}

% 11 --------------------------------------------------------------------------
\begin{frame}{Мини-проверка: нулевое скалярное произведение}
\vspace{0.30em}

\begin{columns}[T,onlytextwidth]
\column{0.46\textwidth}
Пусть
\[
  u=(1,2),
  \qquad
  v=(2,-1).
\]

Тогда
\[
  u^\top v
  =1\cdot2+2\cdot(-1)
  =0.
\]

По формуле
\[
  u^\top v=\norm{u}\norm{v}\cos\theta
\]
для ненулевых $u,v$ получаем $\cos\theta=0$, то есть
\[
  \boxed{\theta=90^\circ.}
\]

\column{0.50\textwidth}
\centering
\includegraphics[width=0.90\linewidth]{orthogonal_vectors.pdf}
\end{columns}

\end{frame}


% 12 --------------------------------------------------------------------------
\begin{frame}{Производная: не формула, а локальная модель}
\small
\vspace{0.10em}

\begin{columns}[T,onlytextwidth]
\column{0.49\textwidth}
\begin{block}{Определение производной}
Если существует предел
\[
  f'(x)=\lim_{h\to0}\frac{f(x+h)-f(x)}{h},
\]
то функция называется дифференцируемой в точке $x$.
\end{block}

Для оптимизации важнее всего эквивалентная локальная запись
\[
  \boxed{f(x+h)=f(x)+f'(x)h+o(|h|).}
\]

То есть около $x$ функция отличается от своей касательной лишь на величину, малую по сравнению с $|h|$.

\column{0.47\textwidth}
\centering
\includegraphics[width=0.96\linewidth]{local_linear_model.pdf}
\end{columns}
\end{frame}

% 13 --------------------------------------------------------------------------
\begin{frame}{Линейная модель функции}
\vspace{0.45em}

Зафиксируем точку $x$ и рассмотрим функцию от малого смещения $h$:
\[
  m_x(h)=f(x)+f'(x)h.
\]

\begin{block}{Интерпретация}
$m_x(h)$ --- простая модель того, чему примерно равно $f(x+h)$.
\end{block}

При $h=0$ она точна:
\[
  m_x(0)=f(x).
\]

А коэффициент при $h$ равен локальному наклону $f'(x)$.

\end{frame}

% 14 --------------------------------------------------------------------------
\begin{frame}{Знак производной подсказывает направление}
\vspace{0.45em}

Если $f'(x)>0$, то для малого $h>0$
\[
 f(x+h)-f(x)\approx f'(x)h>0,
\]
то есть движение вправо увеличивает функцию.

\vspace{0.35em}

Если $f'(x)<0$, ситуация обратная.

\begin{alertblock}{В одном измерении}
\centering
Чтобы локально уменьшить функцию, разумно двигаться \textbf{против знака производной}.
\end{alertblock}

\end{frame}

% 15 --------------------------------------------------------------------------
\begin{frame}{Из локального направления уже видна знакомая формула}
\vspace{0.7em}

В одном измерении естественный шаг имеет вид
\[
  \boxed{x_{\text{new}}=x-\alpha f'(x),\qquad \alpha>0.}
\]

\begin{block}{Пока не алгоритм}
Сегодня нас интересует только происхождение направления $-f'(x)$. Вопрос «какое брать $\alpha$?» мы сознательно откладываем.
\end{block}

\vspace{0.5em}

\centering
\coursebox[3.1cm]{$f'(x)>0$}
\quad $\Rightarrow$ \quad
\coursebox[3.1cm]{двигаемся влево}

\vspace{0.35em}
\coursebox[3.1cm]{$f'(x)<0$}
\quad $\Rightarrow$ \quad
\coursebox[3.1cm]{двигаемся вправо}

\end{frame}

% 16 --------------------------------------------------------------------------
\begin{frame}{Стационарная точка ещё не обязательно минимум}
\vspace{0.25em}

\begin{block}{Определение}
В одномерном случае точка $x^\star$ называется \textbf{стационарной}, если
\[
  f'(x^\star)=0.
\]
\end{block}

Рассмотрим три функции:
\[
  f_1(x)=x^2,
  \qquad
  f_2(x)=-x^2,
  \qquad
  f_3(x)=x^3.
\]
Во всех трёх случаях $x^\star=0$ стационарна, но тип точки различен:
\[
\begin{array}{c|c|c}
 x^2 & -x^2 & x^3\\ \hline
 \text{минимум} & \text{максимум} & \text{ни минимум, ни максимум}
\end{array}
\]

\begin{alertblock}{Вывод}
Первый порядок фиксирует нулевой наклон, но не определяет тип стационарной точки. Нужна информация второго порядка.
\end{alertblock}
\end{frame}


% 17 --------------------------------------------------------------------------
\begin{frame}{Что будем называть локальной кривизной?}
\vspace{0.25em}

В одной переменной $f'(x)$ описывает локальный наклон, а $f''(x)$ --- скорость изменения этого наклона.

\begin{block}{Определение для этой лекции}
\textbf{Локальной кривизной} одномерной функции в точке $x$ будем называть величину
\[
  f''(x).
\]
Для функции многих переменных кривизну будем измерять \textbf{в выбранном направлении} через вторую производную соответствующего одномерного сечения.
\end{block}

Например,
\[
  f_1(x)=x^2,
  \qquad
  f_2(x)=-x^2,
\]
и в точке $0$ обе функции имеют нулевой наклон, но
\[
  f_1''(0)=2>0,
  \qquad
  f_2''(0)=-2<0.
\]
\end{frame}


% 18 --------------------------------------------------------------------------
\begin{frame}{Квадратичная модель видит изгиб функции}
\small
\vspace{0.15em}

\begin{columns}[T,onlytextwidth]
\column{0.47\textwidth}
При малом $h$
\[
  \boxed{
  f(x+h)
  \approx
  f(x)+f'(x)h+\frac12f''(x)h^2.
  }
\]

\begin{block}{Два уровня информации}
$f'(x)h$ задаёт касательную --- локальный наклон.

\vspace{0.25em}
$\frac12f''(x)h^2$ позволяет модели изгибаться и учитывать изменение наклона.
\end{block}

На рисунке линейная и квадратичная модели совпадают с $e^t$ в точке $0$, но второй порядок заметно лучше повторяет изгиб функции рядом с ней.

\column{0.49\textwidth}
\centering
\includegraphics[width=0.98\linewidth]{curvature_taylor_1d.pdf}
\end{columns}

\end{frame}


% 19 --------------------------------------------------------------------------
\begin{frame}{Что заменит $f'$ и $f''$ в нескольких измерениях?}
\vspace{0.8em}

\centering
\coursebox[4.5cm]{Одна переменная\\$f'(x)$}
\qquad
$\Longrightarrow$
\qquad
\coursebox[4.5cm]{Много переменных\\$\ ?\ $}

\vspace{0.8em}

\coursebox[4.5cm]{Одна переменная\\$f''(x)$}
\qquad
$\Longrightarrow$
\qquad
\coursebox[4.5cm]{Много переменных\\$\ ?\ $}

\vspace{0.8em}

\small
Сначала разберём первый порядок. Ответом будет \textbf{градиент}.

\end{frame}

% 20 --------------------------------------------------------------------------
\begin{frame}{Поверхность и линии уровня --- два взгляда на одну функцию}
\small
\vspace{0.20em}

Рассмотрим
\[
  f(x_1,x_2)=x_1^2+2x_2^2.
\]

\begin{columns}[T,onlytextwidth]
\column{0.39\textwidth}
\begin{block}{Поверхность}
Высота над точкой $(x_1,x_2)$ равна значению $f(x_1,x_2)$.
\end{block}

\begin{block}{Линия уровня}
Контур $f(x_1,x_2)=c$ соединяет точки одной высоты. Для этой функции такие контуры --- эллипсы.
\end{block}

По линиям уровня удобно видеть геометрию оптимизации сверху, не рисуя третью координату.

\column{0.57\textwidth}
\centering
\includegraphics[width=0.70\linewidth]{surface_and_contours.pdf}
\end{columns}

\end{frame}


% 21 --------------------------------------------------------------------------
\begin{frame}{Частная производная: меняем одну координату}
\vspace{0.45em}

Для
\[
  f(x_1,\ldots,x_d)
\]
частная производная
\[
  \frac{\partial f}{\partial x_i}(x)
\]
показывает локальное изменение функции, если мы двигаем только координату $x_i$.

\begin{block}{Пример}
Для
\[
 f(x_1,x_2)=x_1^2+2x_2^2
\]
получаем
\[
 \frac{\partial f}{\partial x_1}=2x_1,
 \qquad
 \frac{\partial f}{\partial x_2}=4x_2.
\]
\end{block}

\end{frame}

% 22 --------------------------------------------------------------------------
\begin{frame}{Градиент собирает все частные производные вместе}
\vspace{0.45em}

\begin{block}{Определение}
\[
  \boxed{
  \grad f(x)=
  \begin{pmatrix}
    \dfrac{\partial f}{\partial x_1}(x)\\[0.4em]
    \vdots\\[0.2em]
    \dfrac{\partial f}{\partial x_d}(x)
  \end{pmatrix}
  }
\]
\end{block}

То есть
\[
  \grad f(x)\in\R^d.
\]

\small
Градиент --- не число, а вектор той же размерности, что и параметры $x$.

\end{frame}

% 23 --------------------------------------------------------------------------
\begin{frame}{Первый пример градиента}
\vspace{0.45em}

Для
\[
 f(x_1,x_2)=x_1^2+2x_2^2
\]
имеем
\[
  \boxed{
  \grad f(x_1,x_2)=
  \begin{pmatrix}
    2x_1\\
    4x_2
  \end{pmatrix}.
  }
\]

В точке $(1,1)$:
\[
  \grad f(1,1)=
  \begin{pmatrix}2\\4\end{pmatrix}.
\]

\vspace{0.2em}
\small
Через несколько слайдов мы поймём, что именно геометрически означает эта стрелка $(2,4)$.

\end{frame}

% 24 --------------------------------------------------------------------------
\begin{frame}{Градиент зависит от точки}
\small
\vspace{-0.05em}

\begin{columns}[T,onlytextwidth]
\column{0.43\textwidth}
Для функции
\[
 f(x_1,x_2)=x_1^2+2x_2^2
\]
имеем
\[
  \grad f(x_1,x_2)=(2x_1,4x_2).
\]
Поэтому
\[
  \grad f(1,1)=(2,4),
\]
\[
  \grad f(-1,1)=(-2,4).
\]

\begin{alertblock}{Важно}
Градиент зависит от точки: это \textbf{поле направлений}, а не один фиксированный вектор.
\end{alertblock}

\column{0.53\textwidth}
\centering
\includegraphics[width=0.88\linewidth]{gradient_field_two_points.pdf}
\end{columns}
\end{frame}

\begin{frame}{Дифференцируемость в $\R^d$: формальная локальная модель}
\small
\vspace{0.10em}

\begin{block}{Определение}
Функция $f:\R^d\to\R$ называется \textbf{дифференцируемой в точке $x$}, если существует вектор $g\in\R^d$ такой, что
\[
  \boxed{
  f(x+h)=f(x)+g^\top h+o(\norm{h}),
  \qquad h\to0.
  }
\]
\end{block}

Такой вектор $g$ единственен и равен градиенту:
\[
  g=\grad f(x).
\]

\begin{block}{Почему его компоненты --- частные производные?}
Подставим $h=t e_i$. Тогда
\[
 \frac{f(x+t e_i)-f(x)}{t}=g_i+o(1),
\]
и при $t\to0$ получаем $g_i=\dfrac{\partial f}{\partial x_i}(x)$.
\end{block}
\end{frame}

% 25 --------------------------------------------------------------------------
\begin{frame}{Главная формула первого порядка}
\vspace{0.6em}

Для малого вектора $h$:
\[
  \boxed{
  f(x+h)\approx f(x)+\grad f(x)^\top h.
  }
\]

\begin{block}{Как читать эту формулу}
\begin{itemize}
  \item $h$ говорит, куда мы собираемся двигаться;
  \item $\grad f(x)$ описывает локальный наклон функции;
  \item $\grad f(x)^\top h$ предсказывает изменение $f$.
\end{itemize}
\end{block}

\vspace{0.3em}
\centering
\textbf{Это многомерный аналог } $f(x+h)\approx f(x)+f'(x)h$.

\end{frame}

% 26 --------------------------------------------------------------------------
\begin{frame}{Почему возникает скалярное произведение?}
\vspace{0.5em}

При малом движении $h$ изменение функции примерно равно
\[
  \Delta f\approx \grad f(x)^\top h.
\]

Вспомним:
\[
  \grad f(x)^\top h
  =
  \norm{\grad f(x)}\,\norm{h}\cos\theta.
\]

\begin{alertblock}{Следствие}
\centering
Не только длина шага, но и \textbf{угол между $h$ и градиентом} определяет, растёт функция или убывает.
\end{alertblock}

\end{frame}

% 27 --------------------------------------------------------------------------
\begin{frame}{Направленная производная}
\vspace{0.45em}

Пусть $v$ --- направление. Рассмотрим движение по прямой
\[
  x(t)=x+tv.
\]

\begin{block}{Определение}
Направленная производная --- обычная производная функции вдоль этой прямой:
\[
  D_vf(x)
  =
  \lim_{t\to0}
  \frac{f(x+tv)-f(x)}{t}.
\]
\end{block}

\small
Если $\norm{v}=1$, то $D_vf(x)$ можно интерпретировать как скорость изменения функции на единицу пройденного расстояния.

\end{frame}

% 28 --------------------------------------------------------------------------
\begin{frame}{Направленная производная через градиент}
\small
\vspace{0.10em}

\begin{block}{Утверждение}
Если $f$ дифференцируема в точке $x$, то для любого направления $v$
\[
  \boxed{D_vf(x)=\grad f(x)^\top v.}
\]
\end{block}

\begin{block}{Доказательство}
Из дифференцируемости при $h=tv$:
\[
  f(x+tv)=f(x)+t\,\grad f(x)^\top v+o(|t|).
\]
Следовательно,
\[
  \frac{f(x+tv)-f(x)}{t}
  =\grad f(x)^\top v+o(1).
\]
Переходя к пределу при $t\to0$, получаем требуемую формулу. $\square$
\end{block}
\end{frame}

% 29 --------------------------------------------------------------------------
\begin{frame}{Небольшая задача: изменение в двух направлениях}
\vspace{0.55em}

Пусть в некоторой точке
\[
  \grad f(x)=(3,4).
\]

Что происходит при движении в направлениях
\[
  v_1=(1,0),
  \qquad
  v_2=(0,-1)?
\]

\vspace{0.3em}

\[
  D_{v_1}f(x)=3,
  \qquad
  D_{v_2}f(x)=-4.
\]

\begin{block}{Вывод}
В направлении $v_1$ функция локально растёт, а в направлении $v_2$ --- убывает.
\end{block}

\end{frame}

% 30 --------------------------------------------------------------------------
\begin{frame}{В каком направлении функция растёт быстрее всего?}
\small
\vspace{0.15em}

\begin{block}{Теорема о направлении наискорейшего роста}
Пусть $f$ дифференцируема в точке $x$ и $\grad f(x)\neq0$. Среди всех единичных направлений $\norm{v}=1$
\[
  \max_{\norm{v}=1}D_vf(x)=\norm{\grad f(x)},
\]
причём максимум достигается при
\[
  v_+=\frac{\grad f(x)}{\norm{\grad f(x)}}.
\]
А минимум равен $-\norm{\grad f(x)}$ и достигается при $v_-=-v_+$.
\end{block}

\small
То есть градиент и антиградиент решают две маленькие оптимизационные задачи по выбору направления.
\end{frame}

% 31 --------------------------------------------------------------------------
\begin{frame}{Доказательство теоремы о наискорейшем росте}
\small
\vspace{0.10em}

\begin{block}{Доказательство теоремы}
По формуле направленной производной и неравенству Коши--Буняковского
\[
  D_vf(x)=\grad f(x)^\top v
  \le \norm{\grad f(x)}\,\norm{v}
  =\norm{\grad f(x)}.
\]
При
\[
  v=v_+=\frac{\grad f(x)}{\norm{\grad f(x)}}
\]
получаем равенство:
\[
  D_{v_+}f(x)=\norm{\grad f(x)}.
\]
Для $v_-=-v_+$ аналогично
\[
  D_{v_-}f(x)=-\norm{\grad f(x)}.
\]
Следовательно, $v_+$ и $v_-$ действительно дают максимальный рост и максимальное убывание. $\square$
\end{block}

\small
Само неравенство Коши--Буняковского используем без доказательства.
\end{frame}

% 32 --------------------------------------------------------------------------
\begin{frame}{Градиент --- направление максимального локального роста}
\vspace{0.35em}

Из доказанной теоремы следует
\[
  \boxed{
  \grad f(x)\ \text{задаёт направление максимального локального роста},
  }
\]
а
\[
  \boxed{
  -\grad f(x)\ \text{задаёт направление максимального локального убывания}.
  }
\]

\begin{block}{Важно про слово «локального»}
Теорема сравнивает только мгновенные скорости изменения функции в текущей точке $x$. Она ничего не утверждает о положении глобального минимума.
\end{block}

\small
Если $\grad f(x)=0$, то все направленные производные первого порядка равны нулю, и первого порядка уже недостаточно.
\end{frame}

% 33 --------------------------------------------------------------------------
\begin{frame}{Для уменьшения функции идём против градиента}
\vspace{0.55em}

Если градиент указывает направление самого быстрого роста, то
\[
  -\grad f(x)
\]
указывает направление самого быстрого локального убывания.

\vspace{0.4em}

Отсюда естественным образом появляется формула
\[
  \boxed{x_{k+1}=x_k-\alpha\grad f(x_k).}
\]

\begin{block}{Но сегодня этого недостаточно}
Мы поняли, \textbf{откуда берётся направление}. Позже отдельно разберём, при каких $\alpha$ функция действительно убывает и когда последовательность сходится.
\end{block}

\end{frame}

% 34 --------------------------------------------------------------------------
\begin{frame}{Градиент не «показывает на минимум»}
\small
\vspace{0.10em}

\begin{columns}[T,onlytextwidth]
\column{0.36\textwidth}
Правильное утверждение:
\[
  -\grad f(x)
\]
даёт направление наибольшего \textbf{локального} убывания среди направлений одинаковой длины.

\vspace{0.20em}
Глобальный минимум может находиться совсем в другом направлении: поверхность может иметь изогнутую долину.

\begin{alertblock}{Ключевое слово}
\centering
\textbf{локально}
\end{alertblock}

\column{0.60\textwidth}
\centering
\includegraphics[width=0.98\linewidth]{gradient_not_to_minimum.pdf}
\end{columns}
\end{frame}


% 35 --------------------------------------------------------------------------
\begin{frame}{Линии уровня и градиент}
\small
\vspace{0.15em}

Линия уровня --- множество точек, в которых функция принимает одно значение:
\[
  \{x:\ f(x)=c\}.
\]

Для
\[
  f(x_1,x_2)=x_1^2+2x_2^2
\]
получаем семейство эллипсов.

\begin{columns}[T,onlytextwidth]
\column{0.34\textwidth}
На рисунке отмечены:
\begin{itemize}
  \item несколько уровней $f(x)=c$;
  \item точка $x$;
  \item градиент $\grad f(x)$;
  \item касательное направление к линии уровня.
\end{itemize}

Следующий слайд объяснит, почему эти два направления перпендикулярны.

\column{0.62\textwidth}
\centering
\includegraphics[width=0.80\linewidth]{level_sets_gradient.pdf}
\end{columns}

\end{frame}


% 36 --------------------------------------------------------------------------
\begin{frame}{Почему градиент перпендикулярен линии уровня?}
\small
\vspace{0.08em}

\begin{block}{Утверждение}
Пусть $f:\R^d\to\R$ дифференцируема, а гладкая кривая $\gamma(t)$ лежит на линии уровня $f=c$ и проходит через точку $x$:
\[
  \gamma(0)=x,\qquad f(\gamma(t))\equiv c.
\]
Тогда касательный вектор $v=\gamma'(0)$ ортогонален градиенту:
\[
  \boxed{\grad f(x)^\top v=0.}
\]
\end{block}

\begin{block}{Доказательство}
Так как $f(\gamma(t))$ постоянно, его производная равна нулю. По правилу цепочки
\[
  0=\frac{d}{dt}f(\gamma(t))\bigg|_{t=0}
  =\grad f(\gamma(0))^\top\gamma'(0)
  =\grad f(x)^\top v.
\]
Значит, $\grad f(x)\perp v$. $\square$
\end{block}
\end{frame}

% 37 --------------------------------------------------------------------------
\begin{frame}{Одинаковый минимум --- разная геометрия}
\small
\vspace{0.10em}

Сравним
\[
  f_1(x)=x_1^2+x_2^2,
  \qquad
  f_2(x)=x_1^2+10x_2^2.
\]
Обе функции минимальны в $(0,0)$, но их линии уровня имеют совершенно разную форму.

\begin{columns}[T,onlytextwidth]
\column{0.48\textwidth}
\centering
\includegraphics[width=0.58\linewidth]{contours_round.pdf}

\column{0.48\textwidth}
\centering
\includegraphics[width=0.58\linewidth]{contours_elongated.pdf}
\end{columns}

\vspace{-0.10em}
\begin{alertblock}{Вопрос на будущее}
Почему по узкой вытянутой долине градиентному методу обычно двигаться труднее? Ответ дадут спектр и обусловленность.
\end{alertblock}

\end{frame}


% 38 --------------------------------------------------------------------------
\begin{frame}{От наклона к кривизне в разных направлениях}
\vspace{0.25em}

В $\R^d$ направлений бесконечно много, поэтому одного числа $f''(x)$ уже нет.

\begin{block}{Определение направленной кривизны}
Выберем направление $v$ и одномерное сечение
\[
  \varphi(t)=f(x+tv).
\]
Если $\varphi''(0)$ существует, будем называть
\[
  \boxed{\kappa_v(x):=\varphi''(0)}
\]
\textbf{кривизной функции $f$ в точке $x$ вдоль направления $v$}.
\end{block}

\small
Следующий вопрос: как вычислить $\kappa_v(x)$ сразу по многомерной функции, не выписывая каждое сечение отдельно?
\end{frame}


% 39 --------------------------------------------------------------------------
\begin{frame}{Гессиан}
\vspace{0.35em}

\begin{block}{Определение}
Матрица вторых производных функции $f:\R^d\to\R$ называется гессианом:
\[
\boxed{
\Hess f(x)=
\begin{pmatrix}
\dfrac{\partial^2 f}{\partial x_1^2} & \cdots & \dfrac{\partial^2 f}{\partial x_1\partial x_d}\\[0.8em]
\vdots & \ddots & \vdots\\[0.4em]
\dfrac{\partial^2 f}{\partial x_d\partial x_1} & \cdots & \dfrac{\partial^2 f}{\partial x_d^2}
\end{pmatrix}.
}
\]
\end{block}

\small
Гессиан имеет размер $d\times d$.

\end{frame}

% 40 --------------------------------------------------------------------------
\begin{frame}{Откуда берётся симметрия гессиана?}
\vspace{0.45em}

В позиции $(i,j)$ гессиана стоит смешанная производная
\[
  \frac{\partial^2 f}{\partial x_i\partial x_j},
\]
а в симметричной позиции $(j,i)$ ---
\[
  \frac{\partial^2 f}{\partial x_j\partial x_i}.
\]

Если эти величины совпадают для всех пар $i,j$, то автоматически
\[
  \Hess f(x)=\Hess f(x)^\top.
\]

\begin{alertblock}{Но совпадают ли они всегда?}
Нет без дополнительных предположений. На следующем слайде сформулируем стандартное достаточное условие, которым будем пользоваться в курсе.
\end{alertblock}

\end{frame}


% 41 --------------------------------------------------------------------------

% 41 --------------------------------------------------------------------------
\begin{frame}{Теорема о равенстве смешанных производных}
\small
\vspace{0.20em}

\begin{block}{Теорема Шварца--Клеро}
Пусть функция $f$ имеет вторые частные производные в некоторой окрестности точки $x$, и смешанные производные непрерывны в этой окрестности. Тогда для любых $i,j$
\[
  \boxed{
  \frac{\partial^2 f}{\partial x_i\partial x_j}(x)
  =
  \frac{\partial^2 f}{\partial x_j\partial x_i}(x).
  }
\]
\end{block}

\begin{block}{Удобное следствие}
Если $f\in C^2$ (все частные производные до второго порядка существуют и непрерывны), то её гессиан симметричен во всех точках:
\[
  \boxed{\Hess f(x)=\Hess f(x)^\top.}
\]
\end{block}

\small
Доказательство теоремы в этом курсе не требуется. Для нас важен вывод: у достаточно гладкой целевой функции гессиан --- симметричная матрица, поэтому к нему применима спектральная геометрия.

\end{frame}

\begin{frame}{Пример вычисления гессиана}
\vspace{0.4em}

Пусть
\[
 f(x_1,x_2)=x_1^2+3x_1x_2+2x_2^2.
\]

Сначала
\[
 \grad f(x)=
 \begin{pmatrix}
  2x_1+3x_2\\
  3x_1+4x_2
 \end{pmatrix}.
\]

Теперь дифференцируем компоненты градиента:
\[
 \boxed{
 \Hess f(x)=
 \begin{pmatrix}
  2&3\\
  3&4
 \end{pmatrix}.
 }
\]

\small
Для квадратичной функции гессиан не зависит от точки.

\end{frame}

% 42 --------------------------------------------------------------------------
\begin{frame}{Кривизна вдоль направления = вторая производная сечения}
\small
\vspace{0.08em}

\begin{columns}[T,onlytextwidth]
\column{0.46\textwidth}
\begin{block}{Утверждение}
Пусть $f\in C^2$. Для
\[
  \varphi(t)=f(x+tv)
\]
выполнено
\[
  \boxed{\varphi''(0)=v^\top\Hess f(x)v.}
\]
\end{block}

Именно поэтому квадратичная форма
\[
  v^\top\Hess f(x)v
\]
интерпретируется как локальная кривизна вдоль $v$.

\column{0.50\textwidth}
\centering
\includegraphics[width=0.78\linewidth]{directional_curvature.pdf}
\end{columns}
\end{frame}

\begin{frame}{Доказательство формулы направленной кривизны}
\small
\vspace{0.10em}

Пусть $v=(v_1,\ldots,v_d)^\top$ и $\varphi(t)=f(x+tv)$. По правилу цепочки
\[
  \varphi'(t)
  =\sum_{i=1}^d
  \frac{\partial f}{\partial x_i}(x+tv)\,v_i.
\]
Дифференцируем ещё раз:
\[
  \varphi''(t)
  =\sum_{i=1}^d\sum_{j=1}^d
  \frac{\partial^2 f}{\partial x_j\partial x_i}(x+tv)\,v_jv_i.
\]
Это ровно матричная запись
\[
  \varphi''(t)=v^\top\Hess f(x+tv)v.
\]
Подставляя $t=0$, получаем
\[
  \boxed{\varphi''(0)=v^\top\Hess f(x)v.}\qquad\square
\]
\end{frame}


% 43 --------------------------------------------------------------------------
\begin{frame}{Главная формула второго порядка}
\small
\vspace{0.10em}

\begin{block}{Формула Тейлора второго порядка}
Если $f\in C^2$ в окрестности точки $x$, то при $h\to0$
\[
\boxed{
 f(x+h)
 =f(x)
 +\grad f(x)^\top h
 +\frac12h^\top\Hess f(x)h
 +o(\norm{h}^2).
}
\]
\end{block}

\begin{columns}[T,onlytextwidth]
\column{0.47\textwidth}
\begin{block}{Первый порядок}
$\grad f(x)^\top h$ описывает локальный наклон.
\end{block}

\column{0.47\textwidth}
\begin{block}{Второй порядок}
$\frac12h^\top\Hess f(x)h$ описывает изменение наклона, то есть локальную кривизну.
\end{block}
\end{columns}

\small
Доказательство многомерной формулы Тейлора в этом курсе используем без вывода.
\end{frame}

% 44 --------------------------------------------------------------------------
\begin{frame}{Положительная кривизна: локальная «чаша»}
\vspace{0.30em}

\begin{block}{Определение}
Симметричная матрица $H$ называется \textbf{положительно определённой}, если
\[
  v^\top Hv>0
  \qquad\text{для любого }v\neq0.
\]
Обозначение:
\[
  H\succ0.
\]
\end{block}

Если $H=\Hess f(x)\succ0$, то кривизна функции положительна во всех направлениях: каждое одномерное сечение локально изгибается вверх.

\begin{alertblock}{Геометрический образ}
\centering
Второй порядок вокруг точки имеет форму локальной «чаши».
\end{alertblock}
\end{frame}

% 45 --------------------------------------------------------------------------
\begin{frame}{Три типа локальной кривизны}
\small
\vspace{0.10em}

\begin{columns}[T,onlytextwidth]
\column{0.31\textwidth}
\centering
\textbf{$H\succ0$}

\includegraphics[width=0.96\linewidth]{curvature_bowl.pdf}

Кривизна положительна во всех направлениях: локальная чаша.

\column{0.31\textwidth}
\centering
\textbf{$H\prec0$}

\includegraphics[width=0.96\linewidth]{curvature_dome.pdf}

Кривизна отрицательна во всех направлениях: локальный купол.

\column{0.31\textwidth}
\centering
\textbf{$H$ неопределён}

\includegraphics[width=0.96\linewidth]{curvature_saddle.pdf}

Есть направления обоих знаков: седловая геометрия.
\end{columns}

\vspace{0.25em}
\centering
Знак $v^\top Hv$ отвечает на вопрос, вверх или вниз изгибается одномерное сечение поверхности вдоль конкретного направления $v$.

\end{frame}


% 46 --------------------------------------------------------------------------
\begin{frame}{Что происходит в стационарной точке?}
\vspace{0.45em}

Пусть
\[
  \grad f(x^\star)=0.
\]

Тогда линейная часть исчезает, и локально
\[
  f(x^\star+h)
  \approx
  f(x^\star)
  +\frac12h^\top Hh,
  \qquad
  H=\Hess f(x^\star).
\]

\begin{alertblock}{Значит}
\centering
Чтобы понять тип стационарной точки, надо понять знак квадратичной формы $h^\top Hh$.
\end{alertblock}

Следующий инструмент --- собственные значения симметричной матрицы $H$.

\end{frame}

% 47 --------------------------------------------------------------------------
\begin{frame}{Квадратичная форма на самом простом примере}
\vspace{0.45em}

Пусть
\[
  A=
  \begin{pmatrix}
   2&0\\
   0&5
  \end{pmatrix}.
\]

Тогда
\[
  q(x)=x^\top Ax=2x_1^2+5x_2^2.
\]

\begin{block}{Что видно сразу?}
\begin{itemize}
  \item вдоль оси $x_1$ коэффициент кривизны равен $2$;
  \item вдоль оси $x_2$ коэффициент кривизны равен $5$;
  \item во второй координате функция растёт быстрее.
\end{itemize}
\end{block}

\end{frame}

% 48 --------------------------------------------------------------------------
\begin{frame}{Собственный вектор: направление, которое матрица не поворачивает}
\vspace{0.25em}

\begin{columns}[T,onlytextwidth]
\column{0.53\textwidth}
\begin{block}{Определение}
Ненулевой вектор $v$ называется \textbf{собственным вектором} матрицы $A$, если существует число $\lambda$ такое, что
\[
  \boxed{Av=\lambda v.}
\]
Число $\lambda$ называется соответствующим \textbf{собственным значением}.
\end{block}

\small
Матрица не уводит такой вектор с его прямой: меняется только масштаб, а при $\lambda<0$ ещё и ориентация.

\column{0.43\textwidth}
\centering
\includegraphics[width=1.00\linewidth]{eigenvector_scaling.pdf}
\end{columns}
\end{frame}

% 49 --------------------------------------------------------------------------
\begin{frame}{Диагональная матрица: собственные направления уже видны}
\vspace{0.4em}

Для
\[
  A=\begin{pmatrix}2&0\\0&5\end{pmatrix}
\]
имеем
\[
  A\begin{pmatrix}1\\0\end{pmatrix}
  =2\begin{pmatrix}1\\0\end{pmatrix},
  \qquad
  A\begin{pmatrix}0\\1\end{pmatrix}
  =5\begin{pmatrix}0\\1\end{pmatrix}.
\]

Поэтому собственные значения равны
\[
  \lambda_1=2,
  \qquad
  \lambda_2=5.
\]

\begin{alertblock}{Интерпретация для квадратичной формы}
\centering
Собственные значения измеряют кривизну в главных направлениях.
\end{alertblock}

\end{frame}

% 50 --------------------------------------------------------------------------
\begin{frame}{Почему симметричные матрицы особенно удобны}
\vspace{0.55em}

Для любой вещественной симметричной матрицы $A=A^\top$ существует ортонормированный базис собственных векторов
\[
  v_1,\ldots,v_d.
\]

То есть
\[
  v_i^\top v_j=0\quad(i\ne j),
  \qquad
  \norm{v_i}=1.
\]

\begin{block}{Спектральная теорема --- геометрически}
У симметричной матрицы есть взаимно перпендикулярные \textbf{главные направления}, вдоль которых она действует просто как растяжение на $\lambda_i$.
\end{block}

\small
Доказательство спектральной теоремы в этом курсе нам не понадобится.

\end{frame}

% 51 --------------------------------------------------------------------------
\begin{frame}{Квадратичная форма в собственных координатах}
\vspace{0.45em}

Разложим вектор по ортонормированным собственным векторам:
\[
  x=\sum_{i=1}^d z_i v_i.
\]

Тогда
\[
  \boxed{x^\top Ax=\sum_{i=1}^d \lambda_i z_i^2.}
\]

\begin{block}{Почему это удобно?}
Сложная квадратичная форма после поворота координат превращается в сумму независимых квадратов.
\end{block}

\small
Именно поэтому знаки собственных значений полностью определяют знак квадратичной формы.

\end{frame}

% 52 --------------------------------------------------------------------------
\begin{frame}{Знаки собственных значений и тип кривизны}
\vspace{0.35em}

Для симметричной матрицы $A$:

\begin{columns}[T,onlytextwidth]
\column{0.31\textwidth}
\begin{block}{$\lambda_i>0$ для всех $i$}
\[
 x^\top Ax>0
\]
для $x\ne0$.

$A\succ0$.
\end{block}

\column{0.31\textwidth}
\begin{block}{$\lambda_i<0$ для всех $i$}
\[
 x^\top Ax<0
\]
для $x\ne0$.

$A\prec0$.
\end{block}

\column{0.31\textwidth}
\begin{block}{Есть разные знаки}
В одних направлениях форма положительна, в других отрицательна.

Седловая геометрия.
\end{block}
\end{columns}

\end{frame}

% 53 --------------------------------------------------------------------------
\begin{frame}{Критерий второго порядка для стационарной точки}
\small
\vspace{0.15em}

\begin{block}{Теорема: достаточное условие строгого локального минимума}
Пусть $f\in C^2$ в окрестности $x^\star$, $\grad f(x^\star)=0$ и
\[
  H:=\Hess f(x^\star)\succ0.
\]
Тогда $x^\star$ --- строгий локальный минимум функции $f$.
\end{block}

\begin{block}{Другие знаки гессиана}
Если $H\prec0$, то $x^\star$ --- строгий локальный максимум.
Если у $H$ есть собственные значения разных знаков, то $x^\star$ --- седловая точка.
\end{block}

\begin{alertblock}{Нулевые собственные значения}
Если $H$ лишь полуопределён, критерий второго порядка может не дать ответа.
\end{alertblock}
\end{frame}

\begin{frame}{Доказательство достаточного условия минимума}
\small
\vspace{0.08em}

Пусть $\lambda_{\min}>0$ --- наименьшее собственное значение $H$. Тогда для любого $h$
\[
  h^\top Hh\ge \lambda_{\min}\norm{h}^2.
\]
По формуле Тейлора и условию $\grad f(x^\star)=0$
\[
 f(x^\star+h)-f(x^\star)
 =\frac12h^\top Hh+o(\norm{h}^2).
\]
По определению $o(\norm{h}^2)$ для достаточно малых $h$ имеем
\[
  \left|o(\norm{h}^2)\right|
  \le \frac14\lambda_{\min}\norm{h}^2.
\]
Поэтому при $h\neq0$
\[
 f(x^\star+h)-f(x^\star)
 \ge \frac12\lambda_{\min}\norm{h}^2
      -\frac14\lambda_{\min}\norm{h}^2
 =\frac14\lambda_{\min}\norm{h}^2>0.
\]
Следовательно, $x^\star$ --- строгий локальный минимум. $\square$
\end{frame}

% 54 --------------------------------------------------------------------------
\begin{frame}{Почему две похожие задачи могут быть очень разными}
\vspace{0.35em}

Сравним
\[
 f_1(x)=x_1^2+x_2^2,
 \qquad
 f_2(x)=x_1^2+100x_2^2.
\]

У обеих функций
\[
 \argmin f=\{(0,0)\}.
\]

Но спектры гессианов различаются:
\[
 \Hess f_1=\begin{pmatrix}2&0\\0&2\end{pmatrix},
 \qquad
 \Hess f_2=\begin{pmatrix}2&0\\0&200\end{pmatrix}.
\]

\begin{alertblock}{Вопрос следующей лекции}
Почему разница между маленькой и большой кривизной влияет на трудность оптимизации?
\end{alertblock}

\end{frame}

% 55 --------------------------------------------------------------------------
\begin{frame}{Две локальные модели, которые надо унести с собой}
\vspace{0.45em}

\begin{block}{Первый порядок}
\[
  \boxed{f(x+h)\approx f(x)+\grad f(x)^\top h.}
\]
Градиент говорит, как функция меняется при малом движении.
\end{block}

\vspace{0.35em}

\begin{block}{Второй порядок}
\[
  \boxed{
  f(x+h)\approx
  f(x)+\grad f(x)^\top h+\frac12h^\top\Hess f(x)h.
  }
\]
Гессиан добавляет информацию о кривизне.
\end{block}

\end{frame}

% 56 --------------------------------------------------------------------------
\begin{frame}{Что нужно уметь после сегодняшней лекции}
\vspace{0.25em}

\begin{enumerate}
  \item Вычислить градиент и гессиан простой функции нескольких переменных.
  \item Интерпретировать $\grad f(x)^\top v$ как скорость изменения функции в направлении $v$.
  \item Объяснить, почему $-\grad f(x)$ --- направление максимального локального убывания.
  \item Интерпретировать $v^\top\Hess f(x)v$ как кривизну вдоль направления.
  \item По знакам собственных значений гессиана понять локальную геометрию стационарной точки.
\end{enumerate}

\vspace{0.55em}

\begin{alertblock}{Следующая лекция}
\centering
\textbf{Квадратичные задачи оптимизации и их геометрия}: спектр, обусловленность и почему некоторые долины намного труднее других.
\end{alertblock}

\end{frame}

\end{document}